%! The Complete Solution to Ax = b — Unit 3, Lesson 3.3 — Exit Ticket (Answer Key)
\documentclass[10pt]{article}
\usepackage{linalg-article}
\usepackage{linalg-key}   % pulls in -boxes; do NOT also load -boxes
\newcommand{\TallMath}[1]{$\displaystyle #1\rule[-1.4em]{0pt}{3.2em}$}
\newcommand{\xx}{\mathbf{x}}
\newcommand{\bb}{\mathbf{b}}
\newcommand{\svec}{\mathbf{s}}
\newcommand{\zero}{\mathbf{0}}

\begin{document}

\pageheader{Unit 3, Lesson 3.3}{Exit Ticket --- Answer Key}
\namedateperiod

\vspace{0.15in}

No notes. Show your reasoning.

\begin{enumerate}[leftmargin=1.8em, itemsep=16pt]
  \item \textbf{Complete solution.} For $A=\begin{bmatrix} 1 & 4 \\ 2 & 8 \end{bmatrix}$ and
        $\bb=\begin{bmatrix} 3 \\ 6 \end{bmatrix}$, find a particular solution $\xx_p$, attach the
        nullspace, and write the complete solution. Is the solution set a point, line, or plane?
        \ansline{$R_2-2R_1\Rightarrow\left[\begin{smallmatrix} 1&4&3\\0&0&0 \end{smallmatrix}\right]$: $x_1+4x_2=3$ with $x_2$ free. Set $x_2=0$: $\xx_p=\begin{bsmallmatrix} 3\\0 \end{bsmallmatrix}$; nullspace $\svec=\begin{bsmallmatrix} -4\\1 \end{bsmallmatrix}$. Complete $\xx=\begin{bsmallmatrix} 3\\0 \end{bsmallmatrix}+t\begin{bsmallmatrix} -4\\1 \end{bsmallmatrix}$ --- a \emph{line} through $(3,0)$, off the origin.}
  \item \textbf{Count the solutions.} A matrix $A$ has $4$ columns and rank $r=3$, and $A\xx=\bb$ is
        solvable. How many solutions does it have --- one or infinitely many? Explain using $n-r$.
        \ansline{$n-r=4-3=1$ free variable, so there is a nonzero nullspace (a line of do-nothing vectors). A solvable system then has \emph{infinitely many} solutions: $\xx_p+t\svec$.}
  \item \textbf{Solvable?} Reducing an augmented matrix $[A\mid\bb]$ gives
        $\left[\begin{array}{cc|c} 1 & 2 & 5 \\ 0 & 0 & 3 \end{array}\right]$. Does $A\xx=\bb$ have a
        solution? Explain what the bottom row tells you.
        \ansline{No solution. The bottom row reads $0x_1+0x_2=3$, i.e.\ $0=3$, which is impossible --- the system is inconsistent ($\bb\notin C(A)$).}
\end{enumerate}

\begin{teachernote}
Sort responses into ``got it / stopped at $\xx_p$ / solvability slip.'' Item~1 checks that students
\emph{add the nullspace} (a full-credit answer is the whole line, not just $\xx_p$). Item~3 is the
discriminator: they must read ``$0=$ nonzero'' as \emph{no solution}, not restate the definition. If many
stop at $\xx_p$ on item~1, open 3.4 by re-emphasizing ``particular $+$ the whole nullspace.''
\end{teachernote}

\end{document}
